\documentclass[9pt,twocolumn,twoside]{pnas-new}

\articletype{Biological Sciences}%%%% Article topic/classification

\templatetype{pnasresearcharticle}

\usepackage{xr}
\externaldocument{supp}

\usepackage{flushend} % balance columns on last page

\begin{document}

\title{Seasonal influenza is entrained, not synchronized}

\author[a,1]{Yair Daon}

\affil[a]{Affiliation [TBD --- fill in before submission]}

\leadauthor{Daon}

\significancestatement{US public-health agencies forecast seasonal
  influenza using models that encode connections between states, with
  connection strengths inferred from past mortality data. We show that
  when influenza epidemics rise and fall together across states
  because they share seasonal forcing, the connection strength is
  mathematically impossible to learn from the data: no estimator,
  however sophisticated, can distinguish a connected system from an
  unconnected one. An independent forecast-skill analysis confirms
  that lagged mortality from neighboring states carries no information
  beyond a national-severity signal. Inter-regional coupling should be
  dropped from operational influenza forecasts, simplifying
  preparedness models without loss of accuracy.}

\authorcontributions{Y.D. designed the research, performed the
  analysis, and wrote the paper.}

\authordeclaration{The author declares no competing interest.}

\correspondingauthor{\textsuperscript{1}To whom correspondence should
  be addressed. E-mail: yair.daon@gmail.com}

\keywords{identifiability $|$ Cram\'er-Rao lower bound $|$ seasonal
  influenza $|$ entrainment $|$ epidemic forecasting}

\begin{abstract}
Epidemic models for influenza forecasting routinely encode
inter-regional connectivity, and forecast skill may improve when one
state's lagged incidence is used as a covariate for another's. Whether
this improvement reflects directed inter-regional coupling or some
other mechanism is unclear. We apply the Cramér-Rao Lower Bound (CRLB)
to a two-region epidemic model fitted on 12 seasons of excess
pneumonia and influenza mortality across 47 US states: 95\%
CRLB-derived confidence intervals around the estimated connectivity
overlap zero for nearly all state pairs. Independently, in a SARIMA
model forecast, another state's pneumonia and influenza mortality
provides no improvement in forecast skill beyond the national
average. Connectivity is therefore unsupported by either analysis and
should be dropped from forecast models used for influenza
preparedness, in accordance with Occam's razor---or, more formally, a
conditional Granger causality test.
\end{abstract}

\dates{This manuscript was compiled on \today}

\maketitle
\thispagestyle{firststyle}
\ifthenelse{\boolean{shortarticle}}{\ifthenelse{\boolean{singlecolumn}}{\abscontentformatted}{\abscontent}}{}

\Firstpage

% Introduction: PNAS convention — no "Introduction" heading; suppress.
% intro.tex contains \Endparasplit after the first paragraph to end
% the asymmetric first-page layout.
%% {\renewcommand{\section}[1]{}% === BEGIN intro.tex ===
%% \providecommand{\Endparasplit}{} % no-op for paper.tex; defined in pnas-new.cls

%% \section{Introduction}
Operational influenza forecasting is a cornerstone of US public-health
preparedness. The CDC FluSight collaborative challenge
\citep{reich2019} and its analogues for other pathogens
\citep{johansson2019} produce probabilistic forecasts that inform
vaccine distribution, hospital staffing, and other policy
decisions. It is a consensus in spatial analyses of US influenza and
related respiratory viruses that inter-regional spread is well
described by mobility data---commuting flows from US Census workflow
data~\citep{viboud2006, charu2017}, domestic airline passenger
volumes~\citep{charu2017, zhang2026}, and mobile-phone movement
indices in emergent-outbreak
settings~\citep{kraemer2020}. Consequently, spatial epidemic models
routinely encode inter-regional connectivity as an input derived from
these external sources rather than from case data themselves. A
natural modelling question sits beneath this practice: does
information from other regions actually help forecast a given region's
flu incidence, and if so, what is the source of the improvement?

%% One major characteristic of seasonal influenza epidemics is
%% \emph{entrainment} \citep{pikovsky2001}.  Entrainment occurs when an
%% independent oscillator is driven by a periodic forcing. Entrainment
%% can be indistinguishable from mutual synchronization, but unlike
%% synchronization it requires only a unidirectional coupling. In the
%% United States, all states share similar intrinsic transmission
%% dynamics and are driven by similar annual environmental forcing. Thus,
%% state-level mortality trajectories can be strongly correlated without
%% any inter-regional coupling.  The distinction between synchronization
%% and entrainment leads us to ask: \textbf{are seasonal influenza
%%   epidemics synchronized or only entrained?} An answer to this
%% question has immediate implications for influenza forecasting.
%% Outside epidemiology the same problem arises whenever coupled
%% oscillators are driven by common forcing \citep{pikovsky2001}: in
%% neural circuits \citep{reid2019, friston2011}, ecological
%% metapopulations \citep{bjornstad1999, grenfell1998}, and climate
%% networks \citep{tsonis2004architecture}.

%% In order to distinguish between entrainment and synchronization, we
%% utilize the Cram\'er-Rao Lower Bound (CRLB hereafter)
%% \citep{casella2024statistical, wilke2020}. The CRLB gives a lower
%% bound on the variance of \emph{every} unbiased estimator of a model
%% parameter, and can therefore rule out the existence of any usable
%% inference method for that parameter. This makes the CRLB a
%% \emph{design-stage} diagnostic---applicable to any model parameter and
%% any planned dataset, before any inference procedure is attempted.
%% Lack of practical identifiability has previously been shown to hamper
%% reliable predictions in COVID epidemic
%% models~\citep{gallo2022}. Applying the CRLB to a discrete-time
%% two-region SIR model with seasonal forcing, we show that the
%% connectivity parameter is unidentifiable under entrainment. We then
%% \Endparasplit \noindent employ the CRLB to derive CIs around
%% maximum-likelihood estimates of connectivity using 12 seasons of
%% pneumonia and influenza excess mortality across 47~states. For nearly
%% all state pairs, we cannot distinguish positive from zero
%% connectivity. Building on Occam's razor, we conclude that connectivity
%% should not be included in models of seasonal influenza.

%% We verify the previous analysis by employing a standard model for
%% time-series forecast. We observe an improvement in forecast skill when
%% our forecast model includes information on pneumonia and influenza
%% mortality in another state, see Fig.~\ref{fig:sarima_heatmap}. We then
%% show, by including the national average mortality as a covariate, that
%% the other state's mortality does not improve forecast skill. Forecasts
%% that aggregate national-level information into a ``seasonal severity''
%% signal therefore capture essentially all of the inter-regional signal
%% available in the surveillance data. Again invoking Occam's razor, we
%% arrive at the same conclusion: connectivity should not be included in
%% models of seasonal influenza.

A common claim in the influenza forecasting literature is that
inter-regional coupling improves forecasts. Yang et
al.~\citep{yang2014} established the canonical coupled-metapopulation
approach with US commuting flows, and more recent systems have built
on this idea using mechanistic metapopulations, Bayesian hierarchical
models, and graph or neural-network architectures~\citep{pei2018,
  osthus2021, lu2019, aiken2021, deng2020}.  Standard time-series
methods similarly suggest US-state ILI forecasts benefit from
including other states as covariates~\citep{thivierge2025}, with a
modest but consistent improvement in skill. (For a contrary result in
the European setting, see~\citep{kramer2020}.) We reproduce this
picture for our 47 US states: adding one state's lag-one mortality as
a covariate reduces test-period RMSE for 65\% of target-covariate
pairs (Fig.~\ref{fig:sarima_heatmap}).

\textbf{We challenge this claim in three steps.} First, we show via
simulation that connectivity between regions is fundamentally
unidentifiable when epidemics are entrained by common seasonal
forcing. Second, we apply this analysis to real US data and find that
surveillance records contain no statistical evidence of inter-regional
coupling for any state pair. Third, we return to the forecasting
framework of Fig.~\ref{fig:sarima_heatmap} and show that the apparent
state-covariate signal carries no information beyond a
national-severity covariate.

\Endparasplit \noindent 

\begin{figure}[t!]
\centering
\includegraphics[width=\linewidth]{../pix/sarima_heatmap.pdf}
\caption{Test-period RMSE change when another state's lagged mortality
  is added as a covariate to a forecast of a weekly per-capita
  pneumonia and influenza mortality. Each cell corresponds to a single
  target-covariate pair. Blue indicates the covariate deteriorates the
  forecast, and red indicates the covariate improves the forecast.
  Including another state's mortality as covariate reduces RMSE for
  65\% of target-covariate pair. Values clipped to~$\pm
  10\%$. States are sorted by the median effect. Train on 2010--2019,
  test on 2023--2025.}
\label{fig:sarima_heatmap}
\end{figure}
\section*{Results}

\subsection*{Entrainment Nullifies Power}\label{subsec:crlb_inflated}
If the apparent state-covariate signal of
Fig.~\ref{fig:sarima_heatmap} reflects coupling, then connectivity
(denoted $\theta$) must be inferable from incidence data. We test this
by computing statistical power for $\theta > 0$ in a two-region SIR
model. We compare in-phase and reverse-phase entrainment and assume
optimal uncertainty bounds for $\theta$ via the Cramer-Rao lower bound
(see Methods). Figure~\ref{fig:power} shows this power across a grid
of connectivity values $\theta$ and seasonal forcing amplitudes
(denoted $\delta$), with initial conditions drawn from the empirical
distribution of real-data fits (see Methods and
SI~Fig.~\ref{fig:si_ics}). Under in-phase entrainment, statistical
power is <FILL>, i.e.~no test can reliably detect inter-regional
coupling regardless of its magnitude. Under reverse-phase entrainment,
power is <FILL>. Particularly, the red star marks a strongly connected
flu epidemic ($\theta = 0.05$, $\delta = 0.517$; see
Table~\ref{tab:params} in Methods). For in-phase entrainment, power
at the red star is <FILL>, while for reverse-phase entrainment this
power is <FILL>. Thus, in-phase entrainment makes inter-regional
coupling fundamentally unidentifiable from seasonal influenza
surveillance data.

\begin{figure}[t!]
\centering \includegraphics[width=\linewidth]{../pix/power.pdf}
\caption{Statistical power of a CRLB-based test for $H_1:\theta > 0$
  at $\alpha = 0.05$, with 12-season worth of simulated data.
  \textbf{Top:} in-phase forcing ($\varphi_1 =
  \varphi_2$). \textbf{Bottom:} reverse-phase forcing ($\varphi_1 -
  \varphi_2 = \pi$). Initial conditions are drawn from the empirical
  distribution of per-season, per-state values
  (SI~Fig.~\ref{fig:si_ics}). Red star: parameter choices
  corresponding to a strongly driven flu epdemic ($\theta = 0.05$,
  $\delta = 0.517$).}
\label{fig:power}
\end{figure}

\subsection*{Inter-State Connectivity in the US via Influenza}\label{subsec:us_real_data}
For influenza in the US, where states share strong seasonal forcing
and similar phases, surveillance record should therefore carry no
information about inter-state coupling. We test this prediction
directly by analyzing excess pneumonia and influenza mortality from 47
US states over 12 flu seasons. Figure~\ref{fig:seasons} shows raw
weekly death counts for four representative states, with the flu
seasons used for inference highlighted.
%% ; representative model reconstructions for individual state pairs
%% are shown in the Supplementary Information.

\begin{figure}[t!]
\centering
\includegraphics[width=\linewidth]{../pix/seasons_highlighted.pdf}
\caption{Weekly pneumonia and influenza deaths for four US
  states. Light traces show the full time series; darker segments
  indicate the 20-week flu seasons (starting each November) included
  in the analysis. COVID-19 pandemic years (2020--2022) are
  excluded.}
\label{fig:seasons}
\end{figure}

We estimated $\theta$ via maximum-likelihood and derived corresponding
CRLB for all state pairs across
12~seasons. Figure~\ref{fig:theta_vs_crlb} shows $\hat\theta$ against
$\sqrt{\text{Var}(\hat\theta)}$. The dashed line marks the $\alpha =
0.05$ rejection boundary, hence we cannot reject $H_0:\theta=0$ for
points above this line. For the vast majority of pairs of states the
surveillance data provide insufficient statistical evidence in favor
of $H_1:\theta > 0$. The estimates $\hat\theta$ themselves spread
roughly uniformly over $[0, 0.5]$. The apparent forecast benefit of
Fig.~\ref{fig:sarima_heatmap} therefore cannot be the signature of
directional inter-regional coupling: no estimator can recover such
coupling from this surveillance record.

\begin{figure}[t!]
\centering
\includegraphics[width=\linewidth]{../pix/theta_vs_crlb.pdf}
\caption{Estimated connectivity $\hat\theta$
  vs.\ $\sqrt{\text{Var}(\hat\theta)}$ for US state pairs. The dashed
  line is the $\alpha=0.05$ rejection boundary and we cannot reject
  $H_0:\theta=0$ for points above this line. Black circles:
  non-significant, cannot reject $H_0:\theta=0$. Red stars:
  significant, reject $H_0:\theta=0$ in favor of $H_1:\theta>0$.}
\label{fig:theta_vs_crlb}
\end{figure}



\subsection*{State covariates carry no information beyond the national average}\label{subsec:natavg_results}
We now return to the forecast-skill picture of
Fig.~\ref{fig:sarima_heatmap}. If the apparent benefit of including
another state's lag-one mortality is not directional coupling, then
what is it? Two mechanisms can produce the improvement of the signle
state covariate forecast over the univariate forecast we observed in
Fig.~\ref{fig:sarima_heatmap}: (i) directed predictability, in which
$B$'s past actually forecasts $A$'s future beyond $A$'s own past and
shared seasonality; or (ii) shared national forcing, in which $B$'s
lagged value is a noisy proxy for the current week's national-level
flu-season severity. The two mechanisms make opposite predictions:
%% for the comparison of (state and national average) to (national
%% average)
under (i) we expect a reduction in RMSE going from (national average)
to (state and national average), since $y_{B,t-1}$ carries directional
information that $\bar y_{t-1}$ does not. In contrast, under (ii) we
expect no reduction in RMSE when going from (national average) to
(state and national average), since $\bar y_{t-1}$ carries all the
information pertaining to the severity. Hence, a comparison of RMSEs
between (national average) and (state and national average) will allow
us to discern (i) from (ii).

This comparison is in fact a conditional Granger-causality
test~\citep{granger1969, geweke1984}. We had already established in
Fig.~\ref{fig:sarima_heatmap} that $y_{B,t-1}$ improves forecast of
$y_{A,t}$, so either $y_{B,t-1}\to y_{A,t}$, or there is a common
cause. If the forecast skill using $y_{B,t-1}$ and $\bar{y}_{t-1}$ is
no better than that of $\bar{y}_{t-1}$, then $\bar{y}_{t-1}$ is the
common cause and $Y_B \leftarrow \bar Y \to Y_A$~\citep{pearl2009}. In
this case, the apparent $Y_B \to Y_A$ relationship is confounded by
the shared national-level severity rather than reflecting directional
inter-regional coupling.




\begin{figure}[t!]
\centering
\includegraphics[width=\linewidth]{../pix/sarima_density.pdf}
\caption{Distribution of relative test-period RMSE change
  (\eqref{eq:sarima_delta}) when the covariate state's lag-one
  mortality $y_{B,t-1}$ is added as a second covariate to a SARIMA
  that already includes the lag-one national average $\bar y_{t-1}$
  across target-covariate pairs. Negative values imply that adding the
  covariate state improves the forecast. The distribution is nearly
  symmetric around zero (median $+0.02\%$, dashed red line), showing
  that the covariate state carries no information beyond the national
  average.}
\label{fig:sarima_density}
\end{figure}

The information in $y_{B,t-1}$ is merely a proxy for the
national-level severity, which the national average captures
cleanly. Had connectivity been an important factor, $y_{B,t-1}$ would
have carried information not contained in $\bar y_{t-1}$ and the joint
model would beat the national-average-only model. We see no such
effect. Conditional Granger causality implies the data are consistent
with a common-cause graphical structure $Y_B \leftarrow \bar Y \to
Y_A$, in which national-level severity confounds the seemed
directed-coupling structure $Y_B \to Y_A$.

\section*{Discussion}
We have shown that US influenza surveillance data provide no
statistical evidence for inter-state coupling. Our claim rests on two
classical results: the Cram\'er-Rao Lower Bound from estimation
theory, and conditional Granger causality from time-series
analysis. Specifically, for epidemics entrained by shared seasonal
forcing, the variance of any unbiased estimator of the connectivity is
so large that we cannot determine whether coupling is non-zero. A
maximum-likelihood optimizer may converge and return an estimate
that may appear informative, but confidence intervals around the
estimate overlap zero for nearly all state pairs. Thus, the estimate
is statistically indistinguishable from a null hypothesis of no
coupling. The approximately uniform spread of inferred connectivity is
actually expected.

We complement the CRLB theory with a conditional Granger causality
test: the results of fitting SARIMA models show that, in a forecasting
framework with no mechanistic SIR machinery at all, the covariate
states' lag-one mortality carry no information beyond the lag-one
national average.  For endemic seasonal pathogens such as
influenza~\citep{viboud2006, tamerius2013}, our results therefore
imply a single operational recommendation: \textbf{inter-regional
  coupling is not a necessary component of preparedness forecasting
  models}.

Our CRLB analysis makes a statement about identifiability: no unbiased
estimator of the connectivity parameter can reject the null hypothesis
of zero connectivity in the entrained regime of US seasonal
influenza. Our SARIMA analysis makes it a statement about conditional
Granger causality: we demonstrated that an inter-state covariate is
merely a noisy proxy for the confounder of national-level flu
severity.

%% Our results complement previous studies on \citep{rosenkrantz2022}:
%% any forecast product that depends on an unidentified connectivity
%% parameter inherits that parameter's full uncertainty, so removing
%% the parameter tightens forecast characterisation rather than
%% loosening it.

Recently, \citep{thivierge2025} asked whether ILI forecasting in US
states benefits from including other-state or US-weighted-average ILI
as covariates in various autoregressions. They found modest but
consistent improvements by both RMSE and weighted interval score, and
importantly observed that other-state and US-weighted-average
covariates produce \emph{similar} gains. In our SARIMA framework, the
national-average covariate alone produces a median RMSE reduction of
$7.4\%$ and reduces RMSE in $97\%$ of target-covariate pairs, while
individual states contribute essentially nothing beyond the national
average. Our CRLB explains why: the connectivity parameter is
unidentifiable under entrainment, so the small effect that
\citep{thivierge2025} observed carries very little information about
inter-regional transmission. Our study therefore complements
\citep{thivierge2025} by providing the information-theoretic reason
for some of the results they observe.

Our analysis also carries implications for influenza forecasting in
tropical regions~\citep{tsang2024}. In tropical regions, weaker and
less coherent seasonal forcing produces reverse-phase outbreak
patterns~\citep{tamerius2013}, which our analysis shows may give rise
to the phase asymmetry needed for identifiability. Connectivity
inference from case data may therefore be more feasible in tropical
settings, and may prove useful in forecasting as well. Our CRLB
analysis offers a precise way to characterise this distinction.

\subsection*{Beyond Influenza}
The identifiability check we apply here is not specific to influenza
or to inter-regional coupling. Contemporary spatial and structured
epidemic models routinely carry connectivity matrices, age-specific
contact rates, and more. Because the CRLB bounds the variance of every
unbiased estimator and is computable from the model and a hypothesised
parameter value alone, the same design-stage check can be applied to
any such model before any inference pipeline is committed; parameters
whose CRLB exceeds the scale of their expected effect cannot be
supported by the planned data and should be removed from the model
altogether in accordance with Occam's razor.  Recent work has begun to
characterise identifiability constraints in network-based epidemic
models~\citep{kiss2023parameter, dilauro2020, kiss2024towards},
building on earlier results that transmission rate and network
structure are fundamentally confounded \citep{welch2011}---together
providing a parallel design-stage program for network parameters. The
mathematical structure underlying our particular finding---oscillators
driven by common periodic forcing---also arises outside epidemiology,
in neural circuits receiving shared input~\citep{reid2019,
  pikovsky2001} and in metapopulations driven by environmental
cycles~\citep{bjornstad1999}; the same entrainment-driven
unidentifiability is plausible there, and the same design-stage
diagnostic applies.

\subsection*{Limitations}
Several limitations warrant discussion. First, our model assumes a
single influenza strain with uniform transmission parameters across
all seasons and regions. In reality, $R_0$ varies substantially across
influenza subtypes~\citep{biggerstaff2014}. This variability is
partially absorbed by allowing the initial susceptible fraction
$S_i(0)$ to vary freely across seasons, but the seasonal forcing shape
is fixed.

Second, our observation model uses a Gaussian approximation to the
binomial. Real mortality data may exhibit overdispersion. However, a
negative binomial model would increase the observation variance and
thus may inflate the CRLB even further. Therefore, our Gaussian CRLB
represents an optimistic bound. Any more realistic noise model would
reinforce the conclusion.

Third, we focused on a two-region model with symmetric coupling.
Extensions to asymmetric coupling or larger networks will likely
reveal additional identifiability challenges.

\matmethods{%

\subsection*{Seasonal ARIMA Forecast}\label{subsec:sarimabench}
We fit a seasonal ARIMA model to each state's weekly per-capita
pneumonia and influenza mortality, with and without covariates drawn
from other states. For a target state $A$, the model of a $d$-times
differenced time series is
\begin{equation}\label{eq:sarima_model}
  y_{A,t}  = a_0
    + \sum_{k=1}^{p} \phi_k\, y_{A, t-k}
    + \sum_{j=1}^{q} \psi_j\, \varepsilon_{A, t-j}
    + \Phi\, y_{A, t-52}
    + \boldsymbol{\beta}^{\!\top} \mathbf{x}_{t-1}
    + \varepsilon_{A,t}%,
\end{equation}
which is a SARIMA$(p,d,q)(1,0,0)_{52}$ model with a one-week-lagged
exogenous regressor vector $\mathbf{x}_{t-1}$.
The non-seasonal orders $(p, d, q)$ and the coefficients are chosen on
the training data by minimizing AIC~\citep{hyndman2008}.

For each ordered pair of states $(A, B)$ with $B \neq A$, we
fit~\eqref{eq:sarima_model} under four configurations of
$\mathbf{x}_{t-1}$:
\begin{align*}
  &\text{(univariate)} && \boldsymbol{\beta} = 0; \\
  &\text{(national average)}  && \mathbf{x}_{t-1} = (\bar y_{t-1}); \\
  &\text{(state)}  && \mathbf{x}_{t-1} = (y_{B, t-1}); \\
  &\text{(state and national average)} && \mathbf{x}_{t-1} = (y_{B, t-1},\; \bar y_{t-1});
\end{align*}
where $\bar y_{t} = \sum_{S} y_{S, t} / 47$ is the national
average across all 47 states; in $\mathbf{x}_{t-1}$ we use it at lag one. We train on 2010--2019 and test on
2023--2025, and report the relative change in test-period RMSE,
\begin{equation}\label{eq:sarima_delta}
  \Delta \;=\; \frac{\text{RMSE}_{\mathrm{alt}} - \text{RMSE}_{\mathrm{ref}}}{\text{RMSE}_{\mathrm{ref}}},
\end{equation}
so $\Delta < 0$ means the alternative model improves the forecast over
the reference.

\subsection*{Model}\label{subsec:model}
Consider two interconnected regions with populations $N_1$ and $N_2$.
Let $S_i(t)$, $I_i(t)$, $R_i(t)$ denote the number of susceptible,
infected, and recovered individuals in region $i$ at time $t$, with
$S_i + I_i + R_i = N_i$. Regional coupling is governed by the contact
matrix
%
%
$$
\mathbf{C} =
\begin{bmatrix}
  1 - \theta  &     \theta \\
      \theta  & 1 - \theta
\end{bmatrix}
$$
%
%
where $\theta \in [0, 0.5]$ is the connectivity parameter. The
transmission rate is a standard seasonal forcing ($\delta$ is forcing
amplitude, $\varphi_i$ is a phase offset):
%
%
$$
\beta_i(t) = \beta_0(1 + \delta\sin(2\pi t + \varphi_i)),
$$
%
%
where $t$ is measured in years. Then, the force of infection in region
$i$ is~\citep{edlund2011, keeling2008}
\begin{equation}\label{eq:foi}
\lambda_i(t) = \beta_i(t) \frac{\sum_j C_{ij}\, I_j(t)}{\sum_j C_{ij}\, N_j}.
\end{equation}

The expected incidence is $\mu_i(t) = S_i(t)[1 -
  \exp(-\lambda_i(t))]$, and the recovery probability per time step is
$\exp(-\gamma)$. Time steps are denoted $\Delta t$, and the
discrete-time SIR dynamics are:
\begin{align}\begin{split}\label{eq:discrete_sir}
S_i(t+\Delta t) &= S_i(t) \exp(-\lambda_i(t)) \\
I_i(t+ \Delta t) &= I_i(t) \exp(-\gamma) + \mu_i(t)
\end{split}\end{align}
We use SIR rather than SIRS dynamics because we conduct inference for
each influenza season separately. Model parameters are calibrated to
seasonal influenza following~\citep{shaman2010}; see
Table~\ref{tab:params}.

\begin{table}[h]
\centering
\begin{tabular}{llll}
\toprule
Parameter & Value & Description & Source \textbackslash calc.\\
\midrule
$R^{\text{max}}_0$ & 3.52 & $R_0$ at peak season & \citep{shaman2010} \\
$R^{\text{min}}_0$ & 1.12 & $R_0$ during the summer & \citep{shaman2010} \\
$D$ & 3.24 & Mean infection period (days) & \citep{shaman2010} \\
$\gamma$ & $2.16$ & Weekly recovery rate & $7 / D$ \\
$\beta_{\text{max}}$ & $3.11$ & Peak-season transmission & $R_0^{\text{max}}(1 - e^{-\gamma})$\\
$\beta_{\text{min}}$ & $0.99$ & Low-season transmission & $R_0^{\text{min}}(1 - e^{-\gamma})$ \\
$\beta_0$ & $2.05$ & Mean transmission rate & $(\beta_{\text{min}} + \beta_{\text{max}})/2$ \\
$\delta$ & $0.517$ & Seasonal amplitude & $\beta_{\text{max}} / \beta_0 - 1$\\
$\rho$ & $5.36 \times 10^{-4}$ & Infection Fatality Rate & \citep{rolfes2018}\\
$\varphi$ & & Phase of seasonal forcing & per state \\
\bottomrule
\end{tabular}
\caption{Model parameters calibrated to seasonal influenza with
  $\beta(t) = \beta_0(1+\delta \sin(2\pi t + \varphi))$.}
\label{tab:params}
\end{table}

\subsection*{Observation Model and Likelihood}\label{subsec:noise}
We assume each infection is reported independently with probability
$\rho$, giving rise to binomially distributed observations, which we approximate as Gaussian:
%
%
$$
Y_i(t) = \rho\mu_i(t) + \epsilon_i(t)
$$
%
%
with
%
%
$$
\epsilon_i(t) \sim \mathcal{N}(0, \rho(1{-}\rho)\mu_i(t)).
$$
%
%
%% The Gaussian model is a deliberate simplification. Real
%% surveillance data typically exhibit overdispersion; a negative
%% binomial model with variance $\rho\mu_i + (\rho\mu_i)^2/k$ would be
%% more realistic. At any fixed parameter value, overdispersion
%% increases the observation variance, which reduces the FIM weights
%% and inflates the CRLB. While the MLE itself may shift under a
%% different noise model, the qualitative conclusion---that the CRLB
%% for $\theta$ is large when epidemics are in-phase---is robust,
%% since the underlying mechanism (vanishing sensitivity
%% $\partial\mu/\partial\theta$ under symmetry) is independent of the
%% observation model.

The unknown parameters for a single influenza season are
$\boldsymbol{\phi} = (\theta, S_1(0), I_1(0), S_2(0), I_2(0))$, giving
the log-likelihood:
\begin{equation}\label{eq:loglik}
\ell(\boldsymbol{\phi}) = -\frac{1}{2}\sum_{t=0}^{T-1}\sum_{i=1}^{2}
\left[\log(2\pi \sigma_i^2(t)) + \frac{(Y_i(t) -
    \rho\mu_i(t))^2}{\sigma_i^2(t)}\right]
\end{equation}
where $\sigma_i^2(t) = \rho(1{-}\rho)\mu_i(t)$.

\subsection*{Derivatives}\label{subsec:sens}
Computing the CRLB and optimizing the likelihood both require
partial derivatives of model outputs with respect to parameters. Let $\phi \in
\{\theta, S_1(0), S_2(0), I_1(0), I_2(0)\}$ denote a parameter ($\phi
= \phi_k$ for $k = 1, \dots, 5$). The partial derivative of the force
of infection~\eqref{eq:foi} is:
\begin{equation}\label{eq:sens_foi}
\frac{\partial \lambda_i}{\partial \phi} = \beta_i(t) \left[
\frac{1}{(\mathbf{C}\mathbf{N})_i}\left(
  \frac{\partial \mathbf{C}}{\partial \phi} \mathbf{I}
  + \mathbf{C} \frac{\partial \mathbf{I}}{\partial \phi}\right)_i
- \frac{\lambda_i}{\beta_i (\mathbf{C}\mathbf{N})_i}
  \left(\frac{\partial \mathbf{C}}{\partial \phi} \mathbf{N}\right)_i
\right]
\end{equation}
The contact matrix derivative is

\begin{align*}
\partial \mathbf{C}/\partial \theta =
\begin{bmatrix}
  -1 &  1 \\
  \phantom{-}1 & -1
\end{bmatrix},\\
\qquad
\partial \mathbf{C}/\partial \phi = \mathbf{0} \text{ for } \phi \neq \theta,
\end{align*}
%% so for $\phi\neq\theta$ only the
%% $\mathbf{C}\,\partial\mathbf{I}/\partial\phi$ term in
%% \eqref{eq:sens_foi} survives.
Partial derivatives of the expected incidence are found from the
product rule:
\begin{equation}\label{eq:sens_mu}
  \frac{\partial \mu_i}{\partial \phi} =
  %
  %
  \frac{\partial S_i}{\partial \phi} [1 - e^{-\lambda_i}] + S_i
  e^{-\lambda_i} \frac{\partial \lambda_i}{\partial \phi}.
\end{equation}
Partials of the system state propagate via:
\begin{align}\begin{split}\label{eq:sens}
\frac{\partial S_i(t+\Delta t)}{\partial \phi} &= e^{-\lambda_i}
\frac{\partial S_i}{\partial \phi} - S_i e^{-\lambda_i}
\frac{\partial \lambda_i}{\partial \phi} \\
\frac{\partial I_i(t+\Delta t)}{\partial \phi} &= e^{-\gamma}
\frac{\partial I_i}{\partial \phi} +
\frac{\partial \mu_i}{\partial \phi}
\end{split}\end{align}
with initial conditions $\partial S_i(0)/\partial S_j(0) =
\delta_{ij}$, $\partial I_i(0)/\partial I_j(0) = \delta_{ij}$, where
$\delta_{ij}$ is Kronecker's Delta. We find these partial derivatives
jointly with the state equations. These sensitivity-equation methods
have been reviewed by~\citep{wu2013sensitivity} and recently extended
in~\citep{mester2022}.

\subsection*{Cram\'er-Rao Lower Bound}\label{subsec:crlb}
For a single epidemic season, the Cram\'er-Rao Lower Bound (CRLB)
states that $\text{Var}(\hat\phi_k) \geq [\mathbf{J}^{-1}]_{kk}$,
where $\mathbf{J}$ is the $5 \times 5$ Fisher Information Matrix with
entries $J_{kl} =
\mathbb{E}[(\partial\ell/\partial\phi_k)(\partial\ell/\partial\phi_l)]$.
One can verify that if we define
\begin{equation*}
A_i(t) := -\frac{1}{2\mu_i(t)} + \frac{Y_i(t) -
  \rho\mu_i(t)}{(1{-}\rho)\mu_i(t)} + \frac{(Y_i(t) -
  \rho\mu_i(t))^2}{2\rho(1{-}\rho)\mu_i(t)^2},
\end{equation*}
then differentiating the log-likelihood~\eqref{eq:loglik} yields
%
%
\begin{equation}\label{eq:grad}
\partial \ell / \partial \phi = \sum_{t,i} A_i(t)\,\partial
\mu_i(t)/\partial \phi.
\end{equation}
%
%
Observations are independent across $(i,t)$, and we can easily verify
$\mathbb{E}[A_i(t)] = 0$. We obtain, using moments of the Gaussian
$Y_i(t) -\rho \mu_i(t)$ and~\eqref{eq:grad}:
\begin{align}\label{eq:fim_scalar}
  \begin{split}
    J_{kl} &= \mathbb{E}\left [\frac{\partial \ell}{\partial \phi_k}\frac{\partial \ell}{\partial \phi_l}\right ] \\
    %
    %
    %
    &=\sum_{t,i} \mathbb{E}[A_i^2(t)] \frac{\partial \mu_i(t)}{\partial \phi_k}  \frac{\partial \mu_i(t)}{\partial \phi_l} \\
    %
    %
    %
    &= \frac{1+\rho}{2(1-\rho)}\sum_{t=0}^{T-1}\sum_{i=1}^{2}
    \frac{1}{\mu_i(t)^2} \frac{\partial \mu_i(t)}{\partial \phi_k}
    \frac{\partial \mu_i(t)}{\partial \phi_l}
  \end{split}
\end{align}
%
%
In matrix form, \eqref{eq:fim_scalar} can be written as 
\begin{align}\label{eq:fim}
  \begin{split}
  \mathbf{J} &= \mathbf{G}^T \mathbf{W} \mathbf{G} \\
  &\text{ where } \\
  %
  %
  G_{(i,t),k} &= \partial \mu_i(t)/\partial \phi_k \\
  %
  %
  \mathbf{W} &= \frac{1+\rho}{2(1-\rho)}\,\text{diag}(\mu_1(0)^{-2},
  \ldots, \mu_2(T{-}1)^{-2}).
  \end{split}
\end{align}

When multiple seasons are available, we assume independence across
seasons, treating each season's initial conditions as nuisance
parameters. We aggregate to find the Cram\'er-Rao Lower Bound:
\begin{equation}\label{eq:crlb_multi}
\text{Var}(\hat{\theta}) \geq \left(\sum_{s} \frac{1}{
  [\mathbf{J}^{-1}(s)]_{11}}\right)^{-1},
\end{equation}
where $\mathbf{J}(s)$ is the Fisher information matrix for season $s$.
Seasons where the CRLB computation fails (e.g., due to near-zero
incidence producing a singular Fisher information matrix) contribute
no information and are assigned $[\mathbf{J}^{-1}(s)]_{11} = \infty$.

Under standard regularity conditions, the MLE is asymptotically
efficient: $\hat\theta_{\text{MLE}} \sim \mathcal{N}(\theta,
[\mathbf{J}^{-1}]_{11})$ as the number of observations
grows~\citep{casella2024statistical}. The CRLB therefore also serves
as a practical approximation for the variance of the
maximum-likelihood estimate of $\theta$.

Identifiability in systems-biology models is sometimes explored via
the profile likelihood~\citep{raue2009, shrestha2011} which is defined,
for a single season via
%
%
$$
\ell_{\text{PL}}(\theta) = \max_{S_1(0), S_2(0), I_1(0), I_2(0)} \ell(\theta, S_1(0), S_2(0), I_1(0), I_2(0)).
$$
%
%
The profile likelihood is calculated numerically over a grid of
$\theta$ values, optimizing over all other parameters at each grid
point. Such a calculation would be computationally prohibitive at the
scale of our analysis. The CRLB is effectively a quadratic
approximation to the profile likelihood around the maximum-likelihood
estimate and is considerably cheaper to compute. Identifiability for
ODE-based epidemic models has also been studied via structural
identifiability~\citep{tuncer2018, chowell2023structural} and via
identifiable parameter combinations~\citep{eisenberg2014}; we use the
CRLB because it bounds the \emph{variance} of an estimator under a
hypothesised data design, which is the question we want to answer.

\subsection*{Simulated CRLB Analysis}\label{subsec:sim_crlb}
We compute the per-season CRLB across a grid of connectivity values
$\theta \in [10^{-4}, 10^{-1}]$ and seasonal forcing amplitudes
$\delta \in [0, 1)$, comparing in-phase ($\varphi_1 = \varphi_2$) and
  reverse-phase ($\varphi_1 - \varphi_2 = \pi$) epidemics. At each
  grid point, the initial conditions $(S_i(0)/N_i, I_i(0)/N_i)$ for
  each simulated season and region are drawn independently from the
  empirical distribution of per-season, per-state initial conditions
  fitted to the real data (see SI~Fig.~\ref{fig:si_ics}). We repeat
  each configuration 500 times with fresh IC draws. We then compute
  the statistical power of a one-sided $\alpha = 0.05$ Wald test for
  $H_1:\theta > 0$ under the asymptotic Normal approximation of the
  MLE:
\begin{equation}\label{eq:power}
\text{Power}(\theta, \delta, \varphi) = 1 - \Phi\!\left(z_\alpha -
\frac{\theta}{\sqrt{\text{CRLB}/12}}\right),
\end{equation}
where $P\left(\mathcal{N}(0,1) > z_\alpha\right) = 0.05$, and the
per-season CRLB is divided by 12 to approximate the 12-season
aggregated CRLB of~\eqref{eq:crlb_multi} (exact when per-season
CRLBs are equal).  We report the mean power across the 500 IC
realizations as a function of $(\theta, \delta)$.

\subsection*{Influenza Mortality Data}\label{subsec:real_data}
We analyze pneumonia and influenza mortality across 47 US states.  We
exclude Alaska and Hawaii since we do not have absolute humidity data
for them, and we exclude Arizona since it is a very arid state with
relatively inconsistent absolute humidity values (see Supplementary
Information). We integrate data from CDC FluView
(2010--2018)~\citep{cdcfluview2024} and HealthData.gov
(2017--2025)~\citep{healthdatagov2025}, excluding COVID pandemic years
2020--2022. Since pneumonia and influenza deaths include a year-round
non-influenza baseline, we estimate influenza burden as excess
mortality above a pre-season baseline, computed as the mean weekly
death count during July--October preceding each flu
season~\citep{serfling1963, viboud2006}. The infection fatality rate
$\rho \approx 5.36 \times 10^{-4}$ follows from the first row of Table
1 in~\citep{rolfes2018}. Seasonal phases $\varphi_i$ are estimated
from absolute humidity data~\citep{daon2025, shaman2009} by finding
the phase shift that maximizes the correlation between negative
absolute humidity values and a phase-shifted sine. We define each flu
season as $T = 20$ weekly observations starting
November~1$^{\text{st}}$ and fit a single connectivity $\theta$ shared
across all seasons, with per-season initial conditions $(S_i(0),
I_i(0))$ inferred via maximum-likelihood. We aggregate the per-season
Cram\'er-Rao lower bounds via~\eqref{eq:crlb_multi} to obtain a
bound on the overall estimator variance for the connectivity $\theta$.
 
\subsection*{Maximum-Likelihood Estimation}\label{subsec:mle}
We maximize the log-likelihood (\eqref{eq:loglik}) using
\texttt{SLSQP}~\citep{kraft1988, kraft1994} from
\texttt{NLopt}~\citep{NLopt}. Analytical gradients were calculated
via~\eqref{eq:grad}. To handle potential local minima, we use a
multi-start strategy with 200 random initial points and retain the
solution with the highest likelihood. Pilot experiments show that 20
initial points suffice, so we take an order of magnitude more just to
be safe.
}

\showmatmethods{}

\dataavail{Pneumonia and influenza mortality data are publicly
  available from CDC FluView (2010--2018) and HealthData.gov
  (2017--2025);
  %% these are de-identified, population-level surveillance data and
  %% the study is exempt from IRB review.
  Absolute humidity data underlying phase estimation are
  from~\citep{daon2025}. Code is available at an online
  repository<FILL>.}

\acknow{The author thanks Michael Klots for early contributions to
  this work. Computations were performed on the SLURM cluster at
  Bar-Ilan University.}

\showacknow{}

\bibliography{\detokenize{~/lib}}

\end{document}
